% Part: proof-theory
% Chapter: sequent-calculus
% Section: admissible-derivable

\documentclass[../../../include/open-logic-section]{subfiles}

\begin{document}

\olfileid{pt}{seq}{adm}

\olsection{في القواعد المقبولة والقابلة للاشتقاق}

تتعلق نتائج كثيرة في نظرية البرهان، أو تستعمل نتائج تتعلق، بما إذا كان كل ما
!!{prove} باستعمال قاعدة معينة يجوز أن~!!{prove} من دونها. ومبرهنة حذف
القطع أشهر مثال: فهي تبيّن أن كل متتالية يسعنا أن~!!{prove} باستعمال قواعد
حساب متتاليات مع~\Cut{} يمكننا كذلك أن~!!{prove} بغير~\Cut. ولهذه
الخاصية اسم لما لها من أهمية.

\begin{defn}
تكون القاعدة~$\OLId{latin-looped}{R}$ \emph{مقبولة} في حساب متتاليات إذا أمكن برهنة كل متتالية
يمكن برهنتها في الحساب مع القاعدة~$\OLId{latin-looped}{R}$ بغير استعمال~$\OLId{latin-looped}{R}$ كذلك.
\end{defn}

\begin{ex}\ollabel{ex:land-deriv}
تأخذ القاعدة~$\LeftR{\land}$ صورتين مختلفتين في~\Log{G1c} و\Log{G3c}.
ففي~\Log{G1c} نسختان منها: تضم مقدمة إحداهما الحد الأيسر~$!A$ من
!!{formula} الرئيسة~$!A\land !B$، وتضم مقدمة الأخرى الحد الأيمن~$!B$:
\[
\Axiom$ !A, \OLId{greek-variable}{Gamma} \fCenter \OLId{greek-variable}{Delta}$
\RightLabel{$\LeftR{\land}_\OLMathNumeral{1}$}
\UnaryInf$ !A \land !B, \OLId{greek-variable}{Gamma} \fCenter \OLId{greek-variable}{Delta}$
\DisplayProof
\qquad
\Axiom$!B, \OLId{greek-variable}{Gamma} \fCenter \OLId{greek-variable}{Delta}$
\RightLabel{$\LeftR{\land}_\OLMathNumeral{2}$}
\UnaryInf$!A \land !B, \OLId{greek-variable}{Gamma} \fCenter \OLId{greek-variable}{Delta}$
\DisplayProof
\]
أما~\Log{G3c} فله قاعدة واحدة فقط:
\[
\Axiom$ !A, !B, \OLId{greek-variable}{Gamma} \fCenter \OLId{greek-variable}{Delta}$
\RightLabel{$\LeftR{\land}_\OLMathNumeral{3}$}
\UnaryInf$ !A \land !B, \OLId{greek-variable}{Gamma} \fCenter \OLId{greek-variable}{Delta}$
\DisplayProof
\]
ويعني الوسم~$\LeftR{\star}$ أن الصيغة الرئيسة تقع على اليسار، أي في
المقدّم، ويعني~$\RightR{\star}$ أنها تقع على اليمين، أي في التالي.

يمكننا الآن أن نسأل: ``هل قاعدة~$\LeftR{\land}_\OLMathNumeral{3}$ في~\Log{G3c} مقبولة في
\Log{G1c}؟'' نعم؛ إذ يمكن محاكاة كل استدلال يستعمل~$\LeftR{\land}_\OLMathNumeral{3}$
مباشرةً بقواعد~\Log{G1c} وحدها. وإلى جانب~$\LeftR{\land}_\OLMathNumeral{1}$ و
$\LeftR{\land}_\OLMathNumeral{2}$ سنحتاج كذلك إلى القاعدة البنيوية
$\LeftR{\Contraction}$:
\[
\Axiom$!A, !B, \OLId{greek-variable}{Gamma} \fCenter \OLId{greek-variable}{Delta}$
\RightLabel{$\LeftR{\land}_\OLMathNumeral{1}$}
\UnaryInf$!A \land !B, !B, \OLId{greek-variable}{Gamma} \fCenter \OLId{greek-variable}{Delta}$
\RightLabel{$\LeftR{\land}_\OLMathNumeral{2}$}
\UnaryInf$!A \land !B, !A \land !B, \OLId{greek-variable}{Gamma} \fCenter \OLId{greek-variable}{Delta}$
\RightLabel{\LeftR{\Contraction}}
\UnaryInf$!A \land !B, \OLId{greek-variable}{Gamma} \fCenter \OLId{greek-variable}{Delta}$
\DisplayProof
\]
\end{ex}

محاكاة استدلال يستعمل قاعدةً باستدلالات تستعمل قواعد أخرى إحدى طرق إثبات
قبول القاعدة. لكنها ليست الطريقة الوحيدة؛ فبعض القواعد مقبولة ولا يمكن
محاكاتها هكذا. وتسمى القواعد التي يمكن محاكاتها \emph{قابلة للاشتقاق}.

\begin{defn}
تسمى القاعدة~$\OLId{latin-looped}{R}$،
\[
\AxiomC{$\OLId{latin-looped}{S}_\OLMathNumeral{1}$}
\AxiomC{$\dots$}
\AxiomC{$\OLId{latin-looped}{S}_\OLId{latin-ordinary}{n}$}
\RightLabel{$\OLId{latin-looped}{R}$}
\TrinaryInfC{$\OLId{latin-looped}{S}$}
\DisplayProof
\]
\emph{قابلة للاشتقاق} في حساب إذا وُجد~!!{proof} تخطيطي للمتتالية~$\OLId{latin-looped}{S}$
يستعمل قواعد الحساب وبديهياته، إلى جانب متتاليات ابتدائية إضافية من صور
$\OLId{latin-looped}{S}_\OLMathNumeral{1}$، \dots،~$\OLId{latin-looped}{S}_\OLId{latin-ordinary}{n}$.
\end{defn}

يبيّن~!!{proof} الوارد في~\olref{ex:land-deriv} أن
$\LeftR{\land}_\OLMathNumeral{3}$ قاعدة قابلة للاشتقاق في~\Log{G1c}؛ فهو~!!{proof}
تخطيطي لنتيجة~$\LeftR{\land}_\OLMathNumeral{3}$ من مقدمات~$\LeftR{\land}_\OLMathNumeral{3}$ بوصفها متتاليات ابتدائية،
وهو~!!{proof} لا يستعمل إلا قواعد~\Log{G1c}. (ولم يستعمل أي بديهية من
بديهيات~\Log{G1c}.)

\begin{prop}\ollabel{prop:lor-G3-adm}
قاعدتا~$\LeftR{\land}$ و$\RightR{\lor}$ في~\Log{G3c} قابلتان للاشتقاق
في~\Log{G1c}.
\end{prop}

\begin{proof}
  يرد البرهان الخاص بـ~$\LeftR{\land}$ في~\olref{ex:land-deriv}، ويُترك
  البرهان الخاص بـ~$\RightR{\lor}$ تمرينًا.
\end{proof}

\begin{prob}
بيّن أن قاعدة~$\RightR{\lor}$ في~\Log{G3c} قابلة للاشتقاق في~\Log{G1c}.
\end{prob}

\begin{prob}
لنفرض أن~$\liff$~!!{operator} معرّف، وأن~$!A\liff !B$ اختصار لـ
$(!A\lif !B)\land(!B\lif !A)$. بيّن أن القاعدتين
\begin{enumerate}
\item \Axiom$\OLId{greek-variable}{Gamma}, !A, !B \fCenter \OLId{greek-variable}{Delta}$
\Axiom$\OLId{greek-variable}{Gamma} \fCenter \OLId{greek-variable}{Delta}, !A, !B$
\RightLabel{\LeftR{\liff}}
\BinaryInf$!A \liff !B, \OLId{greek-variable}{Gamma} \fCenter \OLId{greek-variable}{Delta}$
\DisplayProof
\item
\Axiom$!A, \OLId{greek-variable}{Gamma} \fCenter \OLId{greek-variable}{Delta}, !B$
\Axiom$!B, \OLId{greek-variable}{Gamma} \fCenter \OLId{greek-variable}{Delta}, !A$
\RightLabel{\RightR{\liff}}
\BinaryInf$\OLId{greek-variable}{Gamma} \fCenter \OLId{greek-variable}{Delta}, !A \liff !B$
\DisplayProof
\end{enumerate}
قابلتان للاشتقاق في: (أ)~\Log{G1c}، و(ب)~\Log{G3c}.
\end{prob}

\begin{ex}\ollabel{ex:weak-adm}
قاعدة~$\RightR{\Weakening}$ في~\Log{G1c}،
\[
\Axiom$ \OLId{greek-variable}{Gamma} \fCenter \OLId{greek-variable}{Delta}$
\RightLabel{\RightR{\Weakening}}
\UnaryInf$ \OLId{greek-variable}{Gamma} \fCenter \OLId{greek-variable}{Delta}, !A$
\DisplayProof
\]
مقبولة في~\Log{G3c}. وفكرة البرهان بسيطة: افترض وجود~!!a{proof}~$\pi$
للمقدمة~$\OLId{greek-variable}{Gamma}\Sequent\OLId{greek-variable}{Delta}$ في~\Log{G3c}. أضف~!!{formula}~$!A$ إلى
تالي كل متتالية في~$\pi$. تبقى البديهيات بديهيات، وتبقى الاستدلالات الصحيحة
صحيحة. وبعد إضافة~!!{formula}~$!A$ تكون المتتالية النهائية لـ~!!{proof} الجديد هي
$\OLId{greek-variable}{Gamma}\Sequent\OLId{greek-variable}{Delta},!A$.

غير أن قاعدة~$\RightR{\Weakening}$ ليست قابلة للاشتقاق في~\Log{G3c}.
فلو كانت كذلك، لوجد~!!a{proof} لـ~$\OLId{greek-variable}{Gamma}\Sequent\OLId{greek-variable}{Delta},!A$ لا يستعمل إلا
بديهيات~\Log{G3c} وقواعده، وربما المتتالية الابتدائية الإضافية
$\OLId{greek-variable}{Gamma}\Sequent\OLId{greek-variable}{Delta}$. وإذا كان هذا~!!{proof} تخطيطيًا محضًا، كما يقتضي
اشتقاق~$\RightR{\Weakening}$، أمكن تحويله إلى~!!a{proof} لـ~$\Sequent !A$
بحذف~$\OLId{greek-variable}{Gamma}$ و$\OLId{greek-variable}{Delta}$. وهذا يحول المتتاليات الابتدائية الإضافية
$\OLId{greek-variable}{Gamma}\Sequent\OLId{greek-variable}{Delta}$ إلى المتتالية الخالية~$\quad\Sequent\quad$.
ولأن المتتالية الخالية لا تحتوي أي~!!{formula}s، فيما تتطلب كل قواعد
\Log{G3c}~!!{formula} فاعلة واحدة على الأقل في مقدماتها، فلا يجوز أن تكون
المتتالية الخالية مقدمة لأي استدلال في هذا~!!{proof} التخطيطي. ومن ثم لا
يجوز أن يكون~!!{proof} تخطيطيًا إلا إذا لم يستعمل بالفعل المتتالية الإضافية
$\OLId{greek-variable}{Gamma}\Sequent\OLId{greek-variable}{Delta}$ متتاليةً ابتدائية. أي إنه سيكون اشتقاقًا تخطيطيًا
لـ~$\OLId{greek-variable}{Gamma}\Sequent\OLId{greek-variable}{Delta},!A$ باستعمال بديهيات~\Log{G3c} وقواعده وحدها،
وسنكون قد بينا أن لكل متتالية من هذه الصورة~!!a{proof} في~\Log{G3c}. غير أن
هذا محال، لأن~\Log{G3c} سليم، فلا يستطيع أن~!!{prove} إلا متتاليات صادقة
في كل~!!{structure}. وإذا أخذنا~$\OLId{greek-variable}{Gamma}=\OLId{greek-variable}{Delta}=\emptyset$ و
$!A\ident\lfalse$، للزم أن~!!{prove}~\Log{G3c}، مثلًا، المتتالية
$\quad\Sequent\lfalse$، وهو لا يفعل.
\end{ex}

لنقدم برهانًا استقرائيًا أنيقًا لكون الإضعاف قاعدة مقبولة في~\Log{G3c}.
بل تصح نتيجة أقوى قليلًا، كما توضح الحجة الحدسية: فإضافة~$!A$ إلى يمين كل
متتالية في~$\pi$ لا تغيّر بنية~!!{proof}، ولا تزيد ارتفاعه على وجه الخصوص.
ولأهمية ذلك نضع له تعريفًا مستقلًا.

\begin{defn}
تكون القاعدة~$\OLId{latin-looped}{R}$،
\[
\AxiomC{$\OLId{latin-looped}{S}_\OLMathNumeral{1}$}
\AxiomC{$\dots$}
\AxiomC{$\OLId{latin-looped}{S}_\OLId{latin-ordinary}{n}$}
\RightLabel{$\OLId{latin-looped}{R}$}
\TrinaryInfC{$\OLId{latin-looped}{S}$}
\DisplayProof
\]
\emph{مقبولة مع حفظ~!!{height}}، أو \emph{مقبولة مع~!!{hp}}، في حساب إذا
كان لكل براهين~$\pi_\OLId{latin-ordinary}{i}$ لمقدماتها~$\OLId{latin-looped}{S}_\OLId{latin-ordinary}{i}$، حيث
$\OLId{latin-ordinary}{h}_\OLId{latin-ordinary}{i}=\pheight{\pi_\OLId{latin-ordinary}{i}}$، برهان~$\pi$ للنتيجة~$\OLId{latin-looped}{S}$ يحقق
$\pheight{\pi}\le\max(\OLId{latin-ordinary}{h}_\OLMathNumeral{1},\dots,\OLId{latin-ordinary}{h}_\OLId{latin-ordinary}{n})$.
\end{defn}

ويكفي لكي تكون قاعدة مقبولة مع حفظ~!!{height} أنه كلما كان للمقدمات~$\OLId{latin-looped}{S}_\OLId{latin-ordinary}{i}$
!!{proof}s~$\pi_\OLId{latin-ordinary}{i}$ ذات~!!{height}~$\OLId{latin-ordinary}{h}_\OLId{latin-ordinary}{i}=\pheight{\pi_\OLId{latin-ordinary}{i}}$، وُجد
!!a{proof}~$\pi$ للنتيجة لا يزيد~!!{height}ه~$\pheight{\pi}$ على أكبر
$\OLId{latin-ordinary}{h}_\OLId{latin-ordinary}{i}$. وعلى وجه الخصوص، إذا كانت لها مقدمة واحدة~$\OLId{latin-looped}{S}_\OLMathNumeral{1}$، كفى أن يكون
$\pheight{\pi}\le\pheight{\pi_\OLMathNumeral{1}}$. وفي حالتي~$\RightR{\Weakening}$ و
$\LeftR{\Weakening}$ يكون بالفعل
$\pheight{\pi}=\pheight{\pi_\OLMathNumeral{1}}$، غير أن المساواة ليست لازمة.

\begin{prop}\ollabel{prop:weak-G3c-adm}
قاعدتا~$\RightR{\Weakening}$ و$\LeftR{\Weakening}$ في~\Log{G1c} مقبولتان
مع حفظ~!!{height} في~\Log{G3c}.
\end{prop}

\begin{proof}
  علينا أن نبيّن أنه إذا كان
  $\Log{G3c}\Proves\OLId{greek-variable}{Gamma}\Sequent\OLId{greek-variable}{Delta}$ بـ~!!{proof}~$\pi$ ذي
  !!{height}~$\pheight{\pi}=\OLId{latin-ordinary}{n}$، وُجد~!!{proof}~$\pi'$ في~\Log{G3c}
  لـ~$\OLId{greek-variable}{Gamma}\Sequent\OLId{greek-variable}{Delta},!A$ وارتفاعه~$\pheight{\pi'}=\OLId{latin-ordinary}{n}$. نبرهن ذلك
  بالاستقراء في~$\OLId{latin-ordinary}{n}$. إذا كان~$\OLId{latin-ordinary}{n}=\OLMathNumeral{0}$، كانت~$\pi$ بديهية: فإما أن تقع صيغة
  ذرية~$!C$ في~$\OLId{greek-variable}{Gamma}$ و$\OLId{greek-variable}{Delta}$ معًا، وإما أن تقع~$\lfalse$ في~$\OLId{greek-variable}{Gamma}$.
  وفي الحالتين تبقى~$\OLId{greek-variable}{Gamma}\Sequent\OLId{greek-variable}{Delta},!A$ بديهية.

  وإذا كان~$\OLId{latin-ordinary}{n}=\pheight{\pi}>\OLMathNumeral{0}$، كان لـ~!!{proof}~$\pi$ استدلال أخير.
  نميّز الحالات بحسب القاعدة المستعملة. ولنورد مثالًا واحدًا: لنفرض أن
  آخر الاستدلال هو~$\RightR{\land}$. عندئذ
  $\OLId{greek-variable}{Delta}=\OLId{greek-variable}{Delta}',!B\land !C$، وتكون~!!{proof}s الجزئية المباشرة
  $\pi_\OLMathNumeral{1}$ و$\pi_\OLMathNumeral{2}$ للاستدلال الأخير ذات المتتاليتين النهائيتين
  $\OLId{greek-variable}{Gamma}\Sequent\OLId{greek-variable}{Delta}',!B$ و$\OLId{greek-variable}{Gamma}\Sequent\OLId{greek-variable}{Delta}',!C$ على الترتيب.
  أي إن~$\pi$ من الصورة:
  \[
  \AxiomC{}
  \RightLabel{$\pi_\OLMathNumeral{1}$}
  \Deduce$\OLId{greek-variable}{Gamma} \fCenter \OLId{greek-variable}{Delta}', !B$
  \AxiomC{}
  \RightLabel{$\pi_\OLMathNumeral{2}$}
  \Deduce$\OLId{greek-variable}{Gamma} \fCenter \OLId{greek-variable}{Delta}', !C$
  \RightLabel{\RightR{\land}}
  \BinaryInf$\OLId{greek-variable}{Gamma} \fCenter \OLId{greek-variable}{Delta}', !B \land !C$
  \DisplayProof
  \]
  ولدينا كذلك
  $\OLId{latin-ordinary}{n}=\max(\pheight{\pi_\OLMathNumeral{1}},\pheight{\pi_\OLMathNumeral{2}})+\OLMathNumeral{1}$. ولذلك يكون
  $\pheight{\pi_\OLMathNumeral{1}}<\OLId{latin-ordinary}{n}$ و$\pheight{\pi_\OLMathNumeral{2}}<\OLId{latin-ordinary}{n}$، فينطبق فرض الاستقراء: توجد
  !!{proof}s~$\pi_\OLMathNumeral{1}'$ و$\pi_\OLMathNumeral{2}'$ لـ~$\OLId{greek-variable}{Gamma}\Sequent\OLId{greek-variable}{Delta}',!B,!A$ و
  $\OLId{greek-variable}{Gamma}\Sequent\OLId{greek-variable}{Delta}',!C,!A$. وبقاعدة~$\RightR{\land}$ ينتج لنا
  !!a{proof}~$\pi'$ لـ~$\OLId{greek-variable}{Gamma}\Sequent\OLId{greek-variable}{Delta}',!B\land !C,!A$:
  \[
  \AxiomC{}
  \RightLabel{$\pi_\OLMathNumeral{1}'$}
  \Deduce$\OLId{greek-variable}{Gamma} \fCenter \OLId{greek-variable}{Delta}', !B, !A$
  \AxiomC{}
  \RightLabel{$\pi_\OLMathNumeral{2}'$}
  \Deduce$\OLId{greek-variable}{Gamma} \fCenter \OLId{greek-variable}{Delta}', !C, !A$
  \RightLabel{\RightR{\land}}
  \BinaryInf$\OLId{greek-variable}{Gamma} \fCenter \OLId{greek-variable}{Delta}', !B \land !C, !A$
  \DisplayProof
  \]
  ويكون
  $\pheight{\pi'}=\max(\pheight{\pi_\OLMathNumeral{1}'},\pheight{\pi_\OLMathNumeral{2}'})+\OLMathNumeral{1}
  \le\max(\pheight{\pi_\OLMathNumeral{1}},\pheight{\pi_\OLMathNumeral{2}})+\OLMathNumeral{1}=\pheight{\pi}$.
  وتعالج قواعد~\Log{G3c} الأخرى على هذا القياس: نطبق فرض الاستقراء على كل
  برهان جزئي مباشر ثم نعيد تطبيق القاعدة الأخيرة. وينطبق البرهان المتماثل
  على الإضعاف الأيسر بإضافة الصيغة إلى مقدم كل متتالية.
\end{proof}

النتيجة~\olref{prop:weak-G3c-adm} مفيدة جدًا. وسنضع ترميزًا لتطبيقها
المتكرر: إذا كان~$\pi$~!!a{proof} لـ~$\OLId{greek-variable}{Gamma}\Sequent\OLId{greek-variable}{Delta}$، فإن
$\pi[\Pi\Sequent\Lambda]$ نتيجة تطبيق مقبولية الإضعاف الأيسر بكل
!!{formula}s في~$\Pi$، ومقبولية الإضعاف الأيمن بكل~!!{formula}s في
$\Lambda$. وهو~!!a{proof} لـ~$\OLId{greek-variable}{Gamma},\Pi\Sequent\OLId{greek-variable}{Delta},\Lambda$.

\end{document}
